In this chapter, the proximal optimal consensus problem for heterogeneous multi-agent systems is addressed, in particular with Lagrangian dynamics along with the control system design for the proposed algorithm of chapter 3. The control system is design with sliding surface feedback linearization, to ensure robustness and convergence of the proposed variables.

\section{Heterogeneous Multi-Agent Dynamics }

    Consider a group of $N$ heterogeneous agents, which are described by the following dynamics
    
    \begin{equation}\label{eq:4.1.1}
        M_i(p_i)\ddot{p}_i + C_i(p_i,\dot{p}_i)\dot{p}_i + G_i(p_i) + F_i(p_i)\dot{p}_i = u_i + d_i(t)
    \end{equation}
    where $p_i \in \mathbf{R}^{nN}$ is the position vector of the $i-$th agents, $\dot{p}_i \in \mathbf{R}^{nΝ}$ is the velocity vector and $\ddot{p}_i \in \mathrm{R}^{nΝ}$ is the acceleration vector. The matrix $M_i(p_i)$ is the mass matrix, the matrix $C(p_i,\dot{p}_i)$ is the Coriolis - Centrifugal Forces matrix, $G_i(p_i)$ is the Weight Force matrix and $D(p_i)$ is the viscous friction parameter vector. The vector $u_i$ is the torque input inside the system and $d_i(t)$ is the bounded disturbances that occur in the system. In a more compact form, the above system can be described in a state space form as
    
    \begin{equation}\label{eq:4.1.2}
        \left[\begin{matrix}
            \dot{p}_i \\
            \ddot{p}_i
        \end{matrix}\right] = \left[\begin{matrix}
             \dot{p}_i\\
            -M_i^{-1}(p_i)[C_i(p_i,\dot{p}_i)\dot{p}_i + G_i(p_i) + D_i(p_i)\dot{p}_i] 
        \end{matrix}\right] + \left[\begin{matrix}
            0 \\
            M_i^{-1}(p_i)
        \end{matrix}\right]\left( u_i + d_i(t) \right)
    \end{equation}
    where $F_i(p_i,\dot{p}_i)$ and $g_i(p_i)$ are the local continuous Lipschitz vector fields, with the output be $x_i := p_i$. So the first order dynamic equation will be
    \begin{equation}\label{eq:4.1.3}
        \dot{p} = F(p,\dot{p}) + g(p)(u + d(t))
    \end{equation}
    
    After the definition of the dynamics for the agents, it is shown below some basic assumptions for the system. Firstly, the disturbances are uniformly bounded as shown in  A.\ref{asm:5}.
    
    \begin{assumption}\label{asm:5}
        There exists constants $d_i > 0$ such that $\left \lVert d_i(t) \right \rVert \leq d_i$ for every $t \geq 0, i = 1,...,N$.
    \end{assumption}
    
    Below, are shown the basic assumptions for the matrices of the dynamic for every agent.
    
    \begin{assumption}\label{asm:6}
        The matrices $M_i$ are symmetrical and positive defined. There exists positive constants $m_{i2} \geq m_{i1} > 0$ such that $\lambda_{max}(M_i(p_i)) \leq m_{i2}$ and $\lambda_{min}(M_i(p_i)) \geq m_{i1}$, for every $p_i \in \mathbf{R}^n, i = 1,...,N$. Thus,
        \begin{equation*}
            m_{i1}I_n \leq M_i(p_i) \leq m_{i2}I_n
        \end{equation*}
    \end{assumption}
    
    \begin{assumption}\label{asm:7}
        The matrix $N(p_i,\dot{p}_i) = \dot{M}(p_i) - 2C(p_i,\dot{p}_i)$ is skew-symmetric, with
        \begin{equation}\label{eq:4.1.4}
            p_i^TN(p_i,\dot{p}_i)p_i = p_i^T(\dot{M}(p_i) - 2C(p_i,\dot{p}_i)))p_i = 0
        \end{equation}
    \end{assumption}
    
    \begin{assumption}\label{asm:8}
        There exists constants $c_{ij} > 0$, such that $\left\lVert \frac{\partial M_i(p_i)}{\partial p_{ij}}\right\rVert \leq c_{ij}$, for every $p_i=[p_{i1},...,p_{in}]^T \in \mathbf{R}^n,i=1,...,N,j=1,...,n$
    \end{assumption}
    
    \begin{remark}\label{rem:4.1}
        From A.\ref{asm:8} it can be shown that
        \begin{equation}\label{eq:4.1.5}
            \left\lVert \dot{M}_i(p_i)\right \rVert = \left \lVert \sum_{j = 1}^{n} \frac{\partial M_i(p_i)}{\partial p_{ij}} \dot{p}_{ij} \right \rVert \leq \sum_{j = 1}^{n} \left \lVert \frac{\partial M_i(p_i)}{\partial p_{ij}}\right \rVert \left \lvert \dot{p}_{ij} \right \rvert \leq \sum_{j = 1}^{n} c_{ij} \left \lvert \dot{p}_{ij} \right \rvert \leq c_{M_i} \left \lVert \dot{p}_i \right \rVert
        \end{equation}
        where $c_{M_i} > 0$.
    \end{remark}
    
    Now that all assumptions for the system have been defined, the next step is the design and implementation of the control system. The control system is based in the methodology of sliding surface and feedback linearization, \cite{HK2001}. The purpose of this controller is to project the error variables inside a surface manifold in order to have robustness and achieve the necessary conditions for the system. The purpose of feedback linearization is to ensure that the nonlinear dynamics will be cancelled from the controller, in order to have a more linear approach. This robust distributed controller will ensure that the NOCGPROX variables will tend to zero in order for the algorithm to converge and achieve optimal proximal consensus.

\section{Distributed Control System Design for NOCGPROX Variables}

    Let us introduce the auxiliary variables 
    \begin{equation}\label{eq:4.2.1}
        \omega_i := \dot{q}_i + \vartheta_i q_i
    \end{equation}
    where $\vartheta_i$ are static gain parameters. The variables $\omega_i(t)$ create a sliding surface for the NOCGPROX variables, in order to introduce robustness to the system. The next step is to calculate these new variables and insert them inside the dynamics of the agents.
    
    \begin{equation*}
        \omega_i = \dot{q}_i + \vartheta_i \left[x_i - \rho \left(t - \left\lfloor \frac{t}{T}\right\rfloor T \right)s_i \left( \left\lfloor \frac{t}{T}\right\rfloor T \right) - \left[ 1 - \rho \left( t - \left\lfloor \frac{t}{T}\right\rfloor T\right)\right] s_i \left(\left\lfloor \frac{t}{T}\right\rfloor T - T\right)\right]
    \end{equation*}
    Now assume that
    \begin{equation}\label{eq:4.2.2}
        b_i(t) = \rho \left(t - \left\lfloor \frac{t}{T}\right\rfloor T \right)s_i \left( \left\lfloor \frac{t}{T}\right\rfloor T \right) + \left[ 1 - \rho \left( t - \left\lfloor \frac{t}{T}\right\rfloor T\right)\right] s_i \left(\left\lfloor \frac{t}{T}\right\rfloor T - T\right)
    \end{equation}
    then,
    \begin{equation}\label{eq:4.2.3}
        \omega_i = \dot{p}_i + \vartheta_ip_i - \dot{b}_i(t) - \vartheta_ib_i(t)
    \end{equation}
    
    When the algorithm is executed, the function $b_i(t)$ generates new sample measurements of the adjacent agents and smoothly incorporates them into the algorithm. The next step is to check if the derivative of $b_i(t)$ exists and is continuous in time. Firstly, for time $t \neq kT$ the derivate will be
    \begin{equation*}
        \begin{align*}
            \dot{b}_i(t) &= \frac{d}{dt}\rho\left(t - \left\lfloor \frac{t}{T}\right\rfloor T \right)s_i\left(\left\lfloor \frac{t}{T} \right\rfloor T \right) + \rho\left(t - \left\lfloor \frac{t}{T}\right\rfloor T \right)\dot{s}_i\left(\left\lfloor \frac{t}{T} \right\rfloor T \right)\\
            &- \frac{d}{dt}\rho\left(t - \left\lfloor \frac{t}{T}\right\rfloor T \right)s_i\left(\left\lfloor \frac{t}{T} \right\rfloor T - T \right)  - \rho\left(t - \left\lfloor \frac{t}{T}\right\rfloor T \right)\dot{s}_i\left(\left\lfloor \frac{t}{T} \right\rfloor T - T \right)
        \end{align*}
    \end{equation*}
    then, the derivate will be
    \begin{equation}\label{eq:4.2.4}
        \dot{b}_i(t) = \frac{d}{dt}\rho\left(t - \left\lfloor \frac{t}{T}\right\rfloor T \right) \left[s_i\left(\left\lfloor \frac{t}{T} \right\rfloor T \right) - s_i\left(\left\lfloor \frac{t}{T} \right\rfloor T - T \right)\right]
    \end{equation}
    
    For time periods $t = kT$ the derivative will have value
    \begin{equation*}
        \begin{align}
            &\lim_{t \to kT^-} \dot{b}_i(t) = \left\frac{d\rho(t)}{dt}\right|_{t = kT}\left[s_i((k - 1)T) - s_i((k-2)T)\right] = 0\\
            &\lim_{t \to kT^+} \dot{b}_i(t) = \left\frac{d\rho(t)}{dt}\right|_{t = kT}\left[s_i(kT) - s_i((k-1)T)\right] = 0
        \end{align}
    \end{equation*}
    since $\left\frac{d\rho}{dt}\right|_{t=kT} = 0$, it can be seen that the function $\dot{b}_i(t)$ is continuous and smooth in its domain. Furthermore, in the same manner can be proved and for $\ddot{b}_i(t)$. So, for $t \neq kT$ it is
    \begin{equation*}
        \begin{align}
            \ddot{b}_i(t) &= \frac{d^2 \rho}{dt^2} \left(t - \left\lfloor \frac{t}{T}\right\rfloor T\right)\left[s_i\left(\left\lfloor \frac{t}{T} \right\rfloor T \right) - s_i\left(\left\lfloor \frac{t}{T} \right\rfloor T - T \right)\right]\\
            & + \frac{d\rho}{dt}\left(t - \left\lfloor \frac{t}{T}\right\rfloor T\right) \left[\dot{s}_i\left(\left\lfloor \frac{t}{T} \right\rfloor T \right) - \dot{s}_i\left(\left\lfloor \frac{t}{T} \right\rfloor T - T \right)\right]
        \end{align}
    \end{equation*}
    thus, the second derivate of the function will have the form
    \begin{equation}\label{eq:4.2.5}
        \ddot{b}_i(t) = \frac{d^2 \rho}{dt^2}\left(t - \left\lfloor \frac{t}{T}\right\rfloor T\right)\left[s_i\left(\left\lfloor \frac{t}{T} \right\rfloor T \right) - s_i\left(\left\lfloor \frac{t}{T} \right\rfloor T - T \right)\right]
    \end{equation}
    For time periods of $t = kT$ the derivative will have the value
    \begin{equation*}
        \begin{align}
            &\lim_{t \to kT^-} \ddot{b}_i(t) = \left\frac{d^2 \rho(t)}{dt^2}\right|_{t = kT} \left[s_i((k - 1)T) - s_i((k-2)T)\right] = 0\\
            &\lim_{t \to kT^+} \ddot{b}_i(t) = \left\frac{d^2 \rho(t)}{dt^2}\right|_{t = kT}\left[s_i(kT) - s_i((k-1)T)\right] = 0
        \end{align}
    \end{equation*}
    since $\left \frac{d^{r}\rho}{dt^{r}} \right|_{t = kT} = 0$ and the derivative $\ddot{b}_i(t)$ is a well-defined continuous function. As a next step, the derivative of the variable $\omega_i(t)$ is going to be computed as
    \begin{equation}\label{eq:4.2.6}
        \begin{align*}
            \dot{\omega}_i = \ddot{p}_i + \vartheta_i \dot{p}_i - (\ddot{b}_i(t) + \vartheta_i \dot{b}_i(t))
        \end{align*}
    \end{equation}
    Now, the equation (\ref{eq:4.1.1}) will be inserted into equation (\ref{eq:4.2.6}). So,
    \begin{equation}\label{eq:4.2.7}
        \dot{\omega}_i = M_i^{-1}(p_i)\left[ u_i + d_i(t) - C_i(p_i,\dot{p}_i)\dot{p}_i - G_i(p_i) - F_i(p_i)\dot{p}_i\right] + \vartheta_i \dot{p}_i - (\ddot{b}_i(t) + \vartheta_i \dot{b}_i(t))
    \end{equation}
    Now let be introduced a candidate Lyapunov function for every agent $i \in \mathcal{V}$.
    \begin{equation}\label{eq:4.2.8}
        V_i := \frac{1}{2}\omega_i^TM_i(p_i)\omega_i
    \end{equation}
    It can be seen that this function is positive definitive as sum of square terms and radially unbounded. The next step is to check the derivative of this function, in order to ensure boundedness of the signals and square integration of the proposed variables. In addition, the control law must have a sign like form in order to ensure and $\mathcal{L}_1$ of the variables. So, the necessary condition will be met in order to prove square integration of the sampled measurements. First, let the expression $Q_i(p_i,\dot{p}_i) = C_i(p_i,\dot{p}_i)\dot{p}_i + G_i(p_i) + F_i(p_i)\dot{p}_i$. The next step is to derive the equation (\ref{eq:4.2.8}). So,

    \begin{equation*}
        \begin{align*}
            \dot{V}_i &= \frac{1}{2}\dot{\omega}_i^TM_i(p_i)\omega_i + \frac{1}{2}\omega_i^TM_i(p_i)\dot{\omega}_i + \frac{1}{2}\omega_i^T\dot{M}_i(p_i)\omega_i\\
            &= \omega_i^TM_i(p_i)\dot{\omega}_i + \frac{1}{2}\omega_i^T\dot{M}_i(p_i)\omega_i\\
            &=  \omega_i^TM_i(p_i)\left[M_i^{-1}(p_i)(u_i + d_i(t) - Q_i(p_i,\dot{p}_i)) + \vartheta_i\dot{p}_i - (\ddot{b}_i(t) + \vartheta_i \dot{b}_i(t))\right] + \frac{1}{2}\omega_i^T\dot{M}_i(p_i)\omega_i\\
            &= \omega_i^T\left[u_i - Q_i(p_i,\dot{p}_i) + M_i(p_i)\left(\vartheta_i\dot{p}_i - (\ddot{b}_i(t) + \vartheta_i \dot{b}_i(t))\right)\right] + \omega_i^T d_i(t) +\frac{1}{2}\omega_i^T\dot{M}_i(p_i)\omega_i\\
            &\leq \omega_i^T\left[u_i - Q_i(p_i,\dot{p}_i)\right] + \left\lVert \omega_i \right\rVert \left\lVert M_i(p_i) \right\rVert \left \lVert \vartheta_i\dot{p}_i - (\ddot{b}_i(t) + \vartheta_i \dot{b}_i(t)) \right\rVert + \left\lVert \omega_i \right\rVert \left\lVert d_i(t) \right\rVert + \frac{1}{2}\omega_i^T \left\lVert \dot{M}_i(p_i)\right\rVert \omega_i
        \end{align*}
    \end{equation*}
    now, the assumptions A.\ref{asm:5} and A.\ref{asm:6} along with remark \ref{rem:4.1} are applied in the above inequality
    \begin{equation*}
        \begin{align*}
            \dot{V}_i &\leq \omega_i^T\left[u_i - Q_i(p_i,\dot{p}_i)\right] + \left\lVert m_{i2}I_n \right\rVert \left\lVert \omega_i \right\rVert \left \lVert \vartheta_i\dot{p}_i - (\ddot{b}_i(t) + \vartheta_i \dot{b}_i(t)) \right\rVert  + d_i\left\lVert \omega_i \right\rVert  +\frac{1}{2}c_{M_i} \lVert \dot{p}_i \rVert\omega_i^T\omega_i\\
            &= \omega_i^T \left[u_i - Q_i(p_i,\dot{p}_i)\right] + m_{i2} \left\lVert \omega_i \right\rVert \left \lVert \vartheta_i\dot{p}_i - (\ddot{b}_i(t) + \vartheta_i \dot{b}_i(t)) \right\rVert  + d_i\left\lVert \omega_i \right\rVert +\frac{1}{2}c_{M_i} \lVert \dot{p}_i \rVert\left\lVert \omega_i \right\rVert^2
        \end{align*}
    \end{equation*}
    The next step is to ensure boundedness of the Lyapunov function. In order to have stabilization in the system, the control signal is chosen as
    \begin{equation}\label{eq:4.2.9}
        u_i = Q_i(p_i, \dot{p}_i) - k_{i_1}\omega_i - \left(y_i(\omega_i, \dot{p}_i, \dot{b}_i, \ddot{b}_i) + d_i\right) \frac{\omega_i}{\lVert \omega_i \rVert_1 + e^{-\left(\lambda_1 + \lambda_2 y_i\left(\omega_i, \dot{p}_i, \dot{b}_i, \ddot{b}_i \right) \right) t}}
    \end{equation}
    where $k_{i_1} > 0, \lambda_1, \lambda_2 > 0$, and
    \begin{equation}\label{eq:4.2.10}
        y_i(\omega_i, \dot{p}_i, \dot{b}_i, \ddot{b}_i) = m_{i2} \lVert \vartheta_i \dot{p}_i - (\ddot{b}_i + \vartheta_i \dot{b}_i) \rVert + \frac{1}{2}c_{M_i} \lVert \dot{p}_i \rVert \lVert \omega_i \rVert
    \end{equation}
    The control signal (\ref{eq:4.2.9}) guarantees that the sliding manifold $\omega_i$ will be bounded and square integrable. It also is going to be absolute integrable. As a result, the error variable will inherit these properties in order to have stability and convergence of the variables to zero. Thus, proving the necessary condition for achieving optimal consensus between the agents. So, the derivative of Lyapunov will become,
    \begin{equation*}
        \begin{align}
            \dot{V}_i &\leq \omega_i^T \left[Q_i(p_i, \dot{p}_i) - k_{i_1}\omega_i - \left(y_i + d_i\right) \frac{\omega_i}{\lVert \omega_i \rVert + e^{-\left(\lambda_1 + \lambda_2 y_i\right)t}} - Q_i(p_i,\dot{p}_i)\right]\\
            & + m_{i2} \left\lVert \omega_i \right\rVert \left \lVert \vartheta_i\dot{p}_i - (\ddot{b}_i(t) + \vartheta_i \dot{b}_i(t)) \right\rVert  + d_i\left\lVert \omega_i \right\rVert +\frac{1}{2}c_{M_i} \lVert \dot{p}_i \rVert\left\lVert \omega_i \right\rVert^2\\
            &= -k_{i_1} \lVert \omega_i \rVert^2 - (y_i + d_i)\frac{\lVert \omega_i \rVert^2}{\lVert \omega_i \rVert + e^{-\left(\lambda_1 + \lambda_2 y_i\right)t}}  + m_{i2} \left\lVert \omega_i \right\rVert \left \lVert \vartheta_i\dot{p}_i - (\ddot{b}_i(t) + \vartheta_i \dot{b}_i(t)) \right\rVert\\
            &+  d_i\left\lVert \omega_i \right\rVert +\frac{1}{2}c_{M_i} \lVert \dot{p}_i \rVert\left\lVert \omega_i \right\rVert^2
        \end{align}
    \end{equation*}
   Now, the term $\frac{\lVert \omega_i \rVert^2}{\lVert \omega_i \rVert + e^{-\left(\lambda_1 + \lambda_2 y_i\right)t}}$ can be written as,
   \begin{equation*}
        \begin{align}
            \frac{\lVert \omega_i \rVert^2}{\lVert \omega_i \rVert + e^{-\left(\lambda_1 + \lambda_2 y_i\right)t}} &= \frac{\lVert \omega_i \rVert^2 + \lVert \omega_i \rVert e^{-\left(\lambda_1 + \lambda_2 y_i\right)t}}{\lVert \omega_i \rVert + e^{-\left(\lambda_1 + \lambda_2 y_i\right)t}} - \frac{\lVert \omega_i \rVert e^{-\left(\lambda_1 + \lambda_2 y_i\right)t}}{\lVert \omega_i \rVert + e^{-\left(\lambda_1 + \lambda_2 y_i\right)t}}\\
            &= \lVert \omega_i \rVert \left(\frac{\lVert \omega_i \rVert + e^{-\left(\lambda_1 + \lambda_2 y_i\right)t}}{\lVert \omega_i \rVert + e^{-\left(\lambda_1 + \lambda_2 y_i\right)t}}\right) - \frac{\lVert \omega_i \rVert e^{-\left(\lambda_1 + \lambda_2 y_i\right)t}}{\lVert \omega_i \rVert + e^{-\left(\lambda_1 + \lambda_2 y_i\right)t}}\\
            & \leq \lVert \omega_i \rVert - e^{-(\lambda_1 + \lambda_2 y_i)t}
        \end{align}
   \end{equation*}
   So, the rewritten and simplified expression can be applied to the derivative of the Lyapunov candidate function. Thus,
   \begin{equation*}
       \begin{align}
           \dot{V}_i &\leq -k_{i_1} \lVert \omega_i \rVert^2 - (y_i + d_i)\frac{\lVert \omega_i \rVert^2}{\lVert \omega_i \rVert + e^{-\left(\lambda_1 + \lambda_2 y_i\right)t}}  + m_{i2} \left\lVert \omega_i \right\rVert \left \lVert \vartheta_i\dot{p}_i - (\ddot{b}_i(t) + \vartheta_i \dot{b}_i(t)) \right\rVert\\
            &+  d_i\left\lVert \omega_i \right\rVert +\frac{1}{2}c_{M_i} \lVert \dot{p}_i \rVert\left\lVert \omega_i \right\rVert^2\\
            &\leq -k_{i_1} \lVert \omega_i \rVert^2 - (y_i + d_i)\left(\lVert \omega_i \rVert - e^{-(\lambda_1 + \lambda_2 y_i)t}\right)  + m_{i2} \left\lVert \omega_i \right\rVert \left \lVert \vartheta_i\dot{p}_i - (\ddot{b}_i(t) + \vartheta_i \dot{b}_i(t)) \right\rVert\\
            & + d_i\left\lVert \omega_i \right\rVert +\frac{1}{2}c_{M_i} \lVert \dot{p}_i \rVert\left\lVert \omega_i \right\rVert^2\\
            &= -k_{i_1} \lVert \omega_i \rVert^2 - (y_i + d_i)\lVert \omega_i \rVert + (y_i + d_i)e^{-(\lambda_1 + \lambda_2 y_i)t}+ d_i\left\lVert \omega_i \right\rVert +\frac{1}{2}c_{M_i} \lVert \dot{p}_i \rVert\left\lVert \omega_i \right\rVert^2\\ 
            &+ m_{i2} \left\lVert \omega_i \right\rVert \left \lVert \vartheta_i\dot{p}_i - (\ddot{b}_i(t) + \vartheta_i \dot{b}_i(t)) \right\rVert\\
            &= -k_{i_1} \lVert \omega_i \rVert^2 - (y_i + d_i)\lVert \omega_i \rVert + (y_i + d_i)e^{-(\lambda_1 + \lambda_2 y_i)t}\\
            &+ \left(m_{i2}\left \lVert \vartheta_i\dot{p}_i - (\ddot{b}_i(t) + \vartheta_i \dot{b}_i(t)) \right\rVert + \frac{1}{2}c_{M_i} \lVert \dot{p}_i \rVert\left\lVert \omega_i \right\rVert+ d_i \right)\lVert \omega_i \rVert\\
            &= -k_{i_1} \lVert \omega_i \rVert^2 + (y_i + d_i)e^{-(\lambda_1 + \lambda_2 y_i)t}\\
            &\leq -k_{i_1} \lVert \omega_i \rVert^2 + y_ie^{-\lambda_2 y_i t} e^{-\lambda_1 t} + d_ie^{-(\lambda_1 + \lambda_2 y_i)t}
       \end{align}
   \end{equation*}
   From the above analysis, it can be seen that the derivative of Lyapunov function has an upper bound,
    \begin{equation}\label{eq:4.2.11}
        \dot{V}_i \leq -k_{i_1} \lVert \omega_i \rVert^2 + y_ie^{-\lambda_2 y_i t} e^{-\lambda_1 t} + d_ie^{-\lambda_1 t}
    \end{equation}

    From the inequality (\ref{eq:4.2.11}) it can be seen that the derivative $\dot{V}_i$ is bounded, from the exponential terms. It can be seen that the exponential terms are $\mathcal{L}_1$ and converge to zero with time. So, we have that $\limsup_{t \to \infty} \dot{V}_i(t) \leq  -k_{i_1} \lVert \omega_i \rVert^2$. Because the Lyapunov function is positive defined and radially unbounded with its derivative be bounded and negative defined, then the function is decreasing overtime with lower bound $V_{\infty}$, such that $\lim_{t \to \infty} V_i(t) = V_{\infty}$, for every agent $i \in \mathcal{V}$. Thus, it has been proved that $\dot{V}_i \leq 0$. Additionally, it has been proven that the sliding surface variables are bounded, $\omega_i \in \mathcal{L}_{\infty}$. Therefore, it can be proven that the error variables can be bounded. As the function is bounded and the sliding surface variables are also bounded, then from equation (\ref{eq:4.2.1}) can be proved that the NOCGPROX variables are also bounded. So,
    \begin{equation*}
        \begin{align}
            \dot{q}_i &= -\vartheta_i q_i + \omega_i\\
            q_i(t) &= e^{-\vartheta_i t}q_i(0) + \int_{0}^{t} e^{-\vartheta_i (t - s)} \omega_i (s) ds\\
            \lVert q_i(t) \rVert & \leq e^{-\vartheta_i t} \lVert q_i(0) \rVert + \int_{0}^{t} e^{-\vartheta_i (t - s)} \lVert \omega_i (s) \rVert ds
        \end{align}
    \end{equation*}
    hence $\lVert \omega_i(s) \rVert \leq c_{\omega}$ then the above inequality will have form,
    \begin{equation*}
        \begin{align}
            \lVert q_i(t) \rVert & \leq e^{-\vartheta_i t} \lVert q_i(0) \rVert + c_{\omega}\int_{0}^{t} e^{-\vartheta_i (t - s)} ds\\
            &= e^{-\vartheta_i t} \lVert q_i(0) \rVert + \frac{c_{\omega}}{\vartheta_i}(1 - e^{-\vartheta_i t})\\
            &\leq c_q := \max \left\{\lVert q_i(0) \rVert,  \frac{c_{\omega}}{\vartheta_i}\right\}
        \end{align}
    \end{equation*}
    So, the variable $q_i \in \mathcal{L}_{\infty}$. The next step is to seek the boundedness of the derivative. From the equation (\ref{eq:4.2.1}) it is
    \begin{equation*}
        \dot{q}_i = -\vartheta_i q_i + \omega_i \Rightarrow \lVert \dot{q}_i \rVert \leq \vartheta_i \lVert q_i \rVert + \lVert \omega_i \rVert \Rightarrow \lVert \dot{q}_i \rVert \leq c_{\dot{q}} := \vartheta_i c_q + c_{\omega}
    \end{equation*}
    Thus the derivative $\dot{q_i} \in \mathcal{L}_{\infty}$. The term $y_ie^{-\lambda_2 y_i t} e^{-\lambda_1 t} \in \mathcal{L}_1$. So, the variable $y_i \in \mathcal{L}_1$. Furthermore, it is implied that $\omega_i \in \mathcal{L}_1$, which can be used to prove that $q_i \in \mathcal{L}_1$. In addition, this implies that all the other signals are bounded. Thus, it has been proven that $\dot{p}_i, \dot{b}_i, \ddot{b}_i \in \mathcal{L}_{\infty}$. Taking into account the boundedness of these signals, along with the NOCGPROX variables, can be seen that $p_i, b_i \in \mathcal{L}_{\infty}$. Therefore, from $b_i(t)$ definition, $s_i(\lfloor \frac{t}{T} \rfloor T) \in \mathcal{\ell}_{\infty}$, hence $\rho(t)$ is bounded by definition. As a result, it has been proven that the sampled measurements of the neighbor agents are bounded $x_j(kT) \in \mathcal{\ell}_{\infty}$. The next step of the stability analysis is to ensure that the error variables are square integrable, in order to ensure that they converge to zero. As it has been proved an upper bound for the Lyapunov derivative, it can be utilized to ensure the necessary condition. Therefore,
    \begin{equation*}
        \begin{align}
            \dot{V}_i &\leq -k_{i_1}\lVert \omega_i \rVert ^2\\
            \int_{0}^{\infty} \dot{V}_i dt & \leq -\int_{0}^{\infty} k_{i_1}\lVert \omega_i \rVert ^2 dt\\
            V_{\infty} - V(0) &\leq -k_{i_1} \int_{0}^{\infty} \lVert \omega_i \rVert ^2dt\\
            \Rightarrow \int_{0}^{\infty} \lVert \omega_i \rVert ^2dt &\leq \frac{V(0) - V_{\infty}}{k_{i_1}} < \infty
        \end{align}
    \end{equation*}
    So, it has been proved that the variable $\omega_i \in \mathcal{L}_2$. The next step is to prove that the error variables are square integrable. From equation (\ref{eq:4.2.1}) it can be shown that the variables $q_i \in \mathcal{L}_2$. So,
    \begin{equation*}
        \begin{align}
            \dot{q}_i &= -\vartheta_i q_i + \omega_i\\
            q_i(t) &= e^{-\vartheta_i t}q_i(0) + \int_{0}^{t} e^{-\vartheta_i (t - s)} \omega_i (s) ds\\
            \lVert q_i(t) \rVert & \leq e^{-\vartheta_i t} \lVert q_i(0) \rVert + \int_{0}^{t} e^{-\vartheta_i (t - s)} \lVert \omega_i (s) \rVert ds\\
            \lVert q_i(t) \rVert^2 & \leq \left(e^{-\vartheta_i t} \lVert q_i(0) \rVert + \int_{0}^{t} e^{-\vartheta_i (t - s)} \lVert \omega_i (s) \rVert ds\right)^2
        \end{align}
    \end{equation*}
    Using Hölder's inequality in the above expression, it can be shown that
    \begin{equation*}
        \lVert q_i(t) \rVert^2  \leq e^{-2\vartheta_i t} \lVert q_i(0) \rVert^2 + \int_{0}^{t} e^{-2\vartheta_i (t - s)} \lVert \omega_i (s) \rVert^2 ds
    \end{equation*}
    Now integrate both parts of the inequality in the interval $[0,\infty)$ it can be seen that
    \begin{equation*}
        \begin{align}
            \int_{0}^{\infty} \lVert q_i(t) \rVert^2 dt & \leq \int_{0}^{\infty} e^{-2\vartheta_i t} \lVert q_i(0) \rVert^2 dt + \int_{0}^{\infty} \int_{0}^{t} e^{-2\vartheta_i (t - s)} \lVert \omega_i (s) \rVert^2 ds dt\\
            &= \frac{\lVert q_i(0) \rVert^2}{2\vartheta_i} + \int_{0}^{\infty} \int_{0}^{t} e^{-2\vartheta_i (t - s)} \lVert \omega_i (s) \rVert^2 ds dt
        \end{align}
    \end{equation*}
    It has been proven that the sliding surface variables are $\omega_i \in \mathcal{L}_2$. So, utilizing Fubini's theorem, the double integral can be computed. Nonetheless, it must be ensured that both integral spaces are measurable. It can be seen that both independent integrals are bounded, and their respected spaces are measurable over time. So, the variable $t$ will be set in the interval $[s, \infty]$ and variable $s$ in the interval $[0, \infty]$. Thus,
    \begin{equation*}
        \begin{align*}
            \int_{0}^{\infty} \lVert q_i(t) \rVert^2 dt & \leq \frac{\lVert q_i(0) \rVert^2}{2\vartheta_i} + \int_{0}^{\infty} \lVert \omega_i (s) \rVert^2 \left( \int_{s}^{\infty} e^{-2\vartheta_i (t - s)}dt\right) ds\\
            &= \frac{\lVert q_i(0) \rVert^2}{2\vartheta_i} + \frac{1}{2\vartheta_i}\underbrace{\int_{0}^{\infty} \lVert \omega_i (s) \rVert^2 ds}_{<\infty}
        \end{align*}
    \end{equation*}
    Therefore, it has been proven that the variables $q_i \in \mathcal{L}_2$ with bound,
    \begin{equation}\label{eq:4.2.12}
        \int_{0}^{\infty} \lVert q_i(t) \rVert^2 dt \leq \frac{\lVert q_i(0) \rVert^2}{2\vartheta_i} + \frac{1}{2\vartheta_i} \int_{0}^{\infty} \lVert \omega_i (s) \rVert^2 ds < \infty
    \end{equation}
    Since it has been proved that $q_i \in \mathcal{L}_{\infty} \cap \mathcal{L}_2$ and $\dot{q}_i \in \mathcal{L}_{\infty}$, then from Barbal$\acute{a}$t's lemma is shown that the $\lim_{t \to \infty} q_i(t) = 0$. Therefore, the necessary condition for the proposed variables has been achieved. Furthermore, because the error variables tend to zero from the regulation problem, then its samples will also converge to zero as the iteration of the algorithm increases. Thus, $\lim_{\nu \to \infty} q_i^{\nu D} = 0$, and the algorithm will converge to the desired stationary point.
    The last step is to prove that the error variables are also $\mathcal{L}_1$. From equation (\ref{eq:4.2.1}) it can be shown
    \begin{equation*}
        \begin{align}
            \dot{q}_i &= -\vartheta_i q_i + \omega_i\\
            q_i(t) &= e^{-\vartheta_i t}q_i(0) + \int_{0}^{t} e^{-\vartheta_i (t - s)} \omega_i (s) ds\\
            \lVert q_i(t) \rVert & \leq e^{-\vartheta_i t} \lVert q_i(0) \rVert + \int_{0}^{t} e^{-\vartheta_i (t - s)} \lVert \omega_i (s) \rVert ds\\
            \int_{0}^{\infty} \lVert q_i(t) \rVert dt & \leq \int_{0}^{\infty} e^{-\vartheta_i t} \lVert q_i(0) \rVert dt + \int_{0}^{\infty} \int_{0}^{t} e^{-\vartheta_i (t - s)} \lVert \omega_i (s) \rVert ds dt\\
             & \leq \frac{\lVert q_i(0) \rVert}{\vartheta_i} + \int_{0}^{\infty} \int_{0}^{t} e^{-\vartheta_i (t - s)} \lVert \omega_i (s) \rVert ds dt
        \end{align}
    \end{equation*}
    In order to have the $\mathcal{L}_1$ integral of $\omega_i$ in the above the expression, Fubini's theorem will be used to transform the planes of the integrals. Nonetheless, it must be ensured that both integral spaces are measurable. It can be seen that both independent integrals are bounded, and their respected spaces are measurable over time. So, the variable $t$ will be set in the interval $[s, \infty]$ and variable $s$ in the interval $[0, \infty]$. Thus,
    \begin{equation*}
        \begin{align}
            \int_{0}^{\infty} \lVert q_i(t) \rVert dt & \leq  \frac{\lVert q_i(0) \rVert}{\vartheta_i} + \int_{0}^{\infty} \lVert \omega_i (s) \rVert \left( \int_{s}^{\infty} e^{-\vartheta_i (t - s)}dt\right) ds\\
            &= \frac{\lVert q_i(0) \rVert}{\vartheta_i} + \frac{1}{\vartheta_i}\underbrace{\int_{0}^{\infty} \lVert \omega_i (s) \rVert ds}_{<\infty}
        \end{align}
    \end{equation*}
    So, it has been proven that $q_i \in \mathcal{L}_1$, with bound as expressed in expression (\ref{eq:4.2.13})
    \begin{equation}\label{eq:4.2.13}
        \int_{0}^{\infty} \lVert q_i(t) \rVert dt \leq \frac{\lVert q_i(0) \rVert}{\vartheta_i} + \frac{1}{\vartheta_i}\underbrace{\int_{0}^{\infty} \lVert \omega_i (t) \rVert dt}_{<\infty} < \infty
    \end{equation}

    It has been proven that the error variable $q_i \in \mathcal{L}_{\infty}\cap\mathcal{L}_1\cap\mathcal{L}_2$ and its derivative $\dot{q}_i \in \mathcal{L}_{\infty}$. These are the required conditions in order to prove that the sampled error variables are square summable. So, this concept is illustrated by the following Lemma \ref{lem:4.1}.
    
    \begin{lemma}\label{lem:4.1}
        if $q_i(t) \in \mathcal{L}_{\infty} \cap \mathcal{L}_1 \cap \mathcal{L}_2$ and $\dot{q}_i(t) \in \mathcal{L}_{\infty}$ then the samples $q_i^{\nu D} \in \mathcal{\ell}_2$, $\nu \in \mathbf{N}$,i.e.
        \begin{equation}\label{eq:4.2.14}
            \sum_{\nu = 0}^{\infty} \left \lVert q_i^{\nu D} \right \rVert^2 < \infty
        \end{equation}
    \end{lemma}
    
    \begin{proof}
        The proof of this lemma begins in a constructive manner. Let be the below derivative
        \begin{equation*}
            \frac{d}{dt}\left((t - kT)q_i^2(t)\right) = q_i^2(t) + 2(t - kT)\dot{q}_i(t)q_i(t)
        \end{equation*}
        the next step is to integrate the above equation in the interval $(kT, (k+1)T)$.
        \begin{equation*}
            \begin{align*}
                \int_{kT}^{(k+1)T} \frac{d}{dt}\left((t - kT)q_i^2(t)\right) dt &= \int_{kT}^{(k+1)T} \left((q_i^2(t) + 2(t - kT)\dot{q}_i(t)q_i(t)\right) dt\\
                (t - kT)q_i^2(t)\bigg|_{kT}^{(k+1)T} &= \int_{kT}^{(k+1)T} q_i^2(t)dt + 2\int_{kT}^{(k+1)T}(t - kT)\dot{q}_i(t)q_i(t)dt\\
                T\left\lVert q_i((k+1)T)\right\rVert^2 &\leq  \int_{kT}^{(k+1)T} \left\lVert q_i(t)\right\rVert^2dt + 2T\int_{kT}^{(k+1)T}\left\lvert\dot{q}_i(t)q_i(t) \right\rvert dt\\
            \end{align*}
        \end{equation*}
        The final step is to sum the above expression from $k = 0...\infty$. So,
        \begin{equation*}
            \begin{align*}
                \sum_{k = 0}^{\infty} T\left\lVert q_i((k+1)T)\right\rVert^2 &\leq \sum_{k = 0}^{\infty} \int_{kT}^{(k+1)T} \left\lVert q_i(t)\right\rVert^2dt + \sum_{k = 0}^{\infty} 2T\int_{kT}^{(k+1)T}\left\lvert\dot{q}_i(t)q_i(t) \right\rvert dt\\
                T\sum_{k = 0}^{\infty} \left\lVert q_i((k+1)T)\right\rVert^2 &\leq  \int_{0}^{\infty} \left\lVert q_i(t)\right\rVert^2dt + 2T \int_{0}^{\infty}\left\lvert\dot{q}_i(t)q_i(t) \right\rvert dt\\
                \sum_{k = 0}^{\infty} \left\lVert q_i((k+1)T)\right\rVert^2 &\leq  \frac{1}{T}\int_{0}^{\infty} \left\lVert q_i(t)\right\rVert^2dt + 2 \int_{0}^{\infty}\left\lvert\dot{q}_i(t)q_i(t) \right\rvert dt\\
            \end{align*}
        \end{equation*}
        Hence $\dot{q}_i(t) \in \mathcal{L}_{\infty}$ is essentially bounded, there exist a constant $M > 0$ such that $\lvert \dot{q}_i(t) \rvert \leq M$ almost everywhere. So,
        \begin{equation*}
            \begin{align*}
                \sum_{k = 0}^{\infty} \left\lVert q_i((k+1)T)\right\rVert^2 &\leq  \frac{1}{T}\int_{0}^{\infty} \left\lVert q_i(t)\right\rVert^2dt + 2 \int_{0}^{\infty}\left\lvert\dot{q}_i(t)q_i(t) \right\rvert dt\\
                &=  \frac{1}{T}\int_{0}^{\infty} \left\lVert q_i(t)\right\rVert^2dt + 2 \int_{0}^{\infty}\left\lvert\dot{q}_i(t) \right\rvert \left\lvert q_i(t) \right\rvert dt\\
                &\leq \frac{1}{T}\int_{0}^{\infty} \left\lVert q_i(t)\right\rVert^2dt + 2 \int_{0}^{\infty} M\left\lvert q_i(t) \right\rvert dt\\
                &= \frac{1}{T}\underbrace{\int_{0}^{\infty} \left\lVert q_i(t)\right\rVert^2dt}_{< \infty} + 2M \underbrace{\int_{0}^{\infty} \left\lvert q_i(t) \right\rvert dt}_{<\infty}
            \end{align*}
        \end{equation*}
        It is known that $q_i(t) \in \mathcal{L}_{1}\cap\mathcal{L}_{2}$. Thus, the above result is finite and bounded. Subsequently,
        \begin{equation*}
            \sum_{k = 0}^{\infty} \left\lVert q_i((k+1)T)\right\rVert^2 \leq \frac{1}{T}\int_{0}^{\infty} \left\lVert q_i(t)\right\rVert^2dt + 2M \int_{0}^{\infty} \left\lvert q_i(t) \right\rvert dt < \infty
        \end{equation*}
        \begin{equation*}
            \sum_{\nu = 0}^{\infty} \left\lVert q_i^{\nu D}\right\rVert^2 < \infty
        \end{equation*}
        Therefore it has been proven that $q_i^{\nu D} \in \mathcal{\ell}_2$.
    \end{proof}
    \newpage
    As, it has been stated previously, the variables $s_i(kT) \in \mathcal{\ell}_{\infty}$, and in extent $x_i(kT) \in \mathcal{\ell}_{\infty}$. It is secured that $x_j(kT) \in \mathcal{\ell}_{\infty}$. From Lemma \ref{lem:4.1} it has been shown that $q_i(kT) \in \mathcal{\ell}_2$. So, the (\ref{eq:3.2.18}) will decrease with each passing iteration since the error term converges to zero and the sum of all computed samples are finite. Therefore, the primal variable converges to the stationary point, which is KKT. Therefore, from theorem \ref{thr:3.1} it has been proven the agents achieve optimal consensus. From this development it can be observed that also the sampled adjacent measurements are also $x_j(kT) \in \mathcal{\ell}_2$.
    
    As a summary, it has been proven that the NOCGPROX variables are converging to zero from Barbal$\acute{a}$t Lemma and its samples are square integrable.  Therefore, with the proposed algorithm \ref{alg:reg-prox-gpda} and the robust distributed control law (\ref{eq:4.2.9}), optimal consensus will be achieved in a sublinear manner with asymptotic stability of the system. 




    